% AMS arrow fills combine real font glyphs for the additional accent forms.
\makeatletter
\newcommand{\ReferenceDefineAccent}[3]{%
  \expandafter\newcommand\csname#1\endcsname{\mathpalette{#2{#3}}}}
\ReferenceDefineAccent{Overrightarrow}{\overarrow@}{\Rightarrowfill@}
\ReferenceDefineAccent{overleftharpoon}{\overarrow@}{\arrowfill@\leftharpoonup\relbar\relbar}
\ReferenceDefineAccent{overrightharpoon}{\overarrow@}{\arrowfill@\relbar\relbar\rightharpoonup}
\ReferenceDefineAccent{overlinesegment}{\overarrow@}{\arrowfill@\mapstochar\relbar\mapsfromchar}
\ReferenceDefineAccent{underlinesegment}{\underarrow@}{\arrowfill@\mapstochar\relbar\mapsfromchar}
\@ifundefined{ReferenceGroupLeft}{}{%
  \ReferenceDefineAccent{overgroup}{\overarrow@}{\arrowfill@\ReferenceGroupLeft\ReferenceGroupRule\ReferenceGroupRight}
  \ReferenceDefineAccent{undergroup}{\underarrow@}{\arrowfill@\ReferenceUnderGroupLeft\ReferenceUnderGroupRule\ReferenceUnderGroupRight}}
\makeatother
% A width-bearing phantom sizes the tilde while accents places it below the nucleus.
\newcommand{\utilde}{\mathpalette\ReferenceUnderTilde}
\newcommand{\ReferenceUnderTilde}[2]{\underaccent{#1\widetilde{\hphantom{#2}}}{#2}}
